\documentclass[a4paper, 11pt]{article}

\usepackage[english]{babel}
\usepackage[utf8]{inputenc}
\usepackage[T1]{fontenc}
\usepackage[margin=2.2cm]{geometry}
\usepackage{booktabs}
\usepackage{float}
\usepackage{amssymb}
\usepackage{hyperref}
\hypersetup{colorlinks=true, linkcolor=blue, urlcolor=blue}

\title{\textbf{FRACTAL Random Forest on EMR: Analysis of the Scalability Results}}
\author{Riya Pokhrel \& Lucas Torri}
\date{November 13, 2026}

\begin{document}
\maketitle


\noindent\textbf{Setup.} Unless stated otherwise, all runs come from the 32-node cluster (13/11):
random forest with 30 trees and depth 10, executors with 3 cores and 20\,GB, partition multiplier $\times 6$.
Each point is a single run. The $\times 3$/$\times 6$ pairs from 12/11 behave like repeats and differ by
at most 2\,\%, which gives an idea of the run-to-run noise.

\section{Strong scaling (fixed data, baseline $n_0 = 8$ executors)}

Total time, with speedup $S = T_8/T_n$ and efficiency $E = S/(n/8)$ in brackets.

\begin{table}[H]
\centering
\begin{tabular}{rlll}
\toprule
Executors & 1\,\% & 5\,\% & 10\,\% \\
\midrule
8  & 240\,s & 881\,s & 1796\,s \\
20 & 131\,s (1.83, 73\,\%) & 421\,s (2.09, 84\,\%) & 725\,s (2.48, 99\,\%) \\
32 & 110\,s (2.18, 55\,\%) & 297\,s (2.97, 74\,\%) & 513\,s (3.50, 88\,\%) \\
40 & 105\,s (2.29, 46\,\%) & 274\,s (3.22, 64\,\%) & 468\,s (3.84, 77\,\%) \\
60 & 95\,s (2.53, 34\,\%)  & 235\,s (3.75, 50\,\%) & 395\,s (4.55, 61\,\%) \\
\bottomrule
\end{tabular}
\caption{Strong scaling on 32 nodes.}
\end{table}

\begin{itemize}
  \item \textbf{Efficiency grows with problem size (Gustafson).} At 32 executors it goes from
  55\,\% to 74\,\% to 88\,\%. Training alone reaches a speedup of 5.09 at 10\,\% with 60 executors.
  \item \textbf{A fixed part of about 70--100\,s does not parallelise (Amdahl):} the file listing
  (31--52\,s, single-threaded on the driver), the preprocessing at 1\,\% (51\,s $\rightarrow$ 42\,s), and the
  26 training stages of the MLlib random forest, which are always the same whatever the data size.
  At 1\,\% with 32 executors or more, tasks last only 120--180\,ms and the utilisation drops to 70--76\,\%.
  With 60 executors at 1\,\%, the file listing alone is 33\,\% of the runtime.
  \item At 10\,\%, the 8-executor run used 8\,\% more CPU time than the 20-executor run, which makes
  $E \approx 100\,\%$ at 20. Memory pressure with only 8 executors is a likely cause, but a single run
  cannot confirm it.
\end{itemize}

\section{The knee: one executor per node}

The training CPU time is constant as long as each node hosts one executor, and rises as soon as
executors share nodes.

\begin{table}[H]
\centering
\begin{tabular}{lcccc}
\toprule
Executors & 20 & 32 & 40 & 60 \\
\midrule
Hosts used & 20 & 31 & 32 & 32 \\
CPU at 1\,\%  & 1.00 & 1.11 & 1.21 & 1.44 \\
CPU at 5\,\%  & 1.00 & 1.05 & 1.17 & 1.41 \\
CPU at 10\,\% & 1.00 & 1.05 & 1.18 & 1.43 \\
\bottomrule
\end{tabular}
\caption{Training CPU time on 32 nodes, relative to 20 executors.}
\end{table}

A controlled comparison confirms it: same data (5\,\%), same code, same number of executors and
partitions, on two cluster sizes.

\begin{table}[H]
\centering
\begin{tabular}{rcccc}
\toprule
Executors & 24 nodes (hosts) & 32 nodes (hosts) & $\Delta$ time & $\Delta$ CPU \\
\midrule
20 & 353\,s (20) & 348\,s (20) & $-1$\,\%  & $-1$\,\% \\
32 & 270\,s (24) & 234\,s (31) & $-13$\,\% & $-12$\,\% \\
40 & 244\,s (24) & 212\,s (32) & $-13$\,\% & $-12$\,\% \\
\bottomrule
\end{tabular}
\caption{Training time at 5\,\% on 24 and 32 nodes.}
\end{table}

Each node has 8 vCPUs, which are 4 physical cores with hyper-threading. A second executor on the same
node runs its tasks on hyper-threads, so each task needs about 40\,\% more CPU time. Going from 32 to
60 executors gives 1.875$\times$ more cores but only 1.16$\times$ (1\,\%), 1.26$\times$ (5\,\%) and
1.30$\times$ (10\,\%) faster runs. This limit comes from the hardware, not from the Spark configuration.

\section{Size-up (fixed executors, more data)}

\begin{table}[H]
\centering
\begin{tabular}{rccc}
\toprule
Fraction & 40 ex.: total / train & 60 ex.: total / train & Train CPU per 1\,\% (60 ex.) \\
\midrule
1\,\%  & 105 / 59\,s   & 95 / 51\,s   & 5.8k CPU$\cdot$s \\
5\,\%  & 274 / 212\,s  & 235 / 176\,s & 5.0k CPU$\cdot$s \\
10\,\% & 468 / 386\,s  & 395 / 318\,s & 4.7k CPU$\cdot$s \\
25\,\% & 1030 / 894\,s & 831 / 708\,s & 4.3k CPU$\cdot$s \\
\bottomrule
\end{tabular}
\caption{Size-up on 32 nodes.}
\end{table}

The runtime grows less than linearly with the data: 25$\times$ more data takes only 8.7$\times$ more
time at 60 executors. Once the fixed costs are absorbed, the growth is close to linear: from 10\,\% to
25\,\% (2.5$\times$ the data) the time grows 2.1--2.3$\times$. The CPU time per row even decreases
slightly with more data. Our explanation, which we did not measure directly: at each tree level, every
task builds a split-statistics histogram whose size depends on the level, not on the number of rows.
Bigger partitions spread that fixed cost over more rows.

\section{Weak scaling (fixed data per executor)}

\begin{table}[H]
\centering
\begin{tabular}{llccc}
\toprule
Data per executor & Run & Total & Train & Weak efficiency (total) \\
\midrule
0.25\,\%  & 1\,\% on 4 executors$^{*}$ & 379\,s & 311\,s & -- \\
0.25\,\%  & 5\,\% on 20 executors  & 421\,s & 348\,s & 90\,\% \\
0.25\,\%  & 10\,\% on 40 executors & 468\,s & 386\,s & 81\,\% \\
\midrule
0.125\,\% & 1\,\% on 8 executors   & 240\,s & 185\,s & -- \\
0.125\,\% & 5\,\% on 40 executors  & 274\,s & 212\,s & 88\,\% \\
\bottomrule
\end{tabular}
\caption{Weak scaling. $^{*}$From the 24-node run with $\times 3$: the placement is identical (one
executor per node) and $\times 3 = \times 6$, so it is comparable.}
\end{table}

Two effects explain the loss: more partitions mean more split statistics to merge on the driver at
every tree level, and at 40 executors, 8 nodes host two executors (Section~2).

\section{Configuration check (course slide 36)}

From 5\,\% of the data upwards, the utilisation is 83--95\,\%, GC takes 2--4\,\% of the task time,
CPU time / run time is 0.95--0.97, there is no spill, the task skew (max/median) is 1.5--2.5 and the
scheduler delay is below 1\,\%. The job is CPU-bound and well balanced, and none of the
misconfiguration symptoms of the slide appear. The partition multiplier is not a lever: $\times 3$ and
$\times 6$ differ by at most 2\,\%. Only the 1\,\% runs with many executors show low utilisation and
high skew (4.5--5.6), because there is too little data per core, not because of a configuration error.

\section{Model quality}

\begin{itemize}
  \item \textbf{Scaling does not change the result.} The mIoU stays between 0.375 and 0.391 and the
  overall accuracy (OA) between 0.839 and 0.846 in all runs. At a given fraction, the mIoU varies by at
  most 0.005 across executor counts; the small differences come from the bootstrap sampling of the
  random forest, which is seeded per partition.
  \item \textbf{More data does not improve the model.} The mIoU goes from 0.388 at 1\,\% to 0.376 at
  25\,\%. Part of this is because the test subset grows with the fraction, so the test sets are not the
  same. Mainly, 30 trees of depth 10 (at most 1024 leaves each) limit the capacity of the model, not the
  amount of data.
  \item \textbf{Class imbalance dominates the mIoU.} Ground and vegetation make up 94\,\% of the test
  points. Bridge and structure get an IoU of 0 and ``other'' at most 0.02, so 3 of the 7 classes pull the
  mean down.
  \item \textbf{Height is the most important feature} (0.28--0.30), ahead of NDVI and Red (0.11--0.22)
  and NumberOfReturns (about 0.1). The height normalisation is the most useful preprocessing step.
  \item \textbf{Validation vs test} (1\,\%, 60 executors): the OA is close (0.836 vs 0.846) but the mIoU
  is not (0.263 vs 0.388). Almost all of the gap comes from water (IoU 0.04 vs 0.74) and building
  (0.31 vs 0.48). With only 1\,\% of the files, the rare classes appear in a few patches, so the mIoU
  depends heavily on the sample. The OA is the stable metric.
\end{itemize}

\section{Extrapolation to 50\,\% and 100\,\% (not measured)}

We extrapolated linearly from the 25\,\% run with 60 executors: the training costs about 4.3k CPU$\cdot$s
per 1\,\% of data and the effective throughput is about 150 CPU$\cdot$s per second. This gives about
45--50\,min of training and \textbf{about 1\,h in total for 100\,\%}, and about 30\,min for 50\,\%. This
assumes that the cached training set still fits in executor memory. If it does not, the
\texttt{MEMORY\_AND\_DISK} storage level falls back to the local disks and the run will be slower, but
it should not fail.

\section{Conclusions}

\begin{enumerate}
  \item The pipeline scales well. At 10\,\% it reaches a speedup of 3.5 at 32 executors ($E = 88\,\%$)
  and 4.55 at 60. The weak-scaling efficiency stays at 81--90\,\%.
  \item Larger datasets scale better, because a fixed overhead of about 70--100\,s does not parallelise:
  the file listing on the driver, the short tasks and the 26 training passes of the random forest.
  \item The main limit is the hardware. Beyond one executor per node, executors share 4 physical cores
  through hyper-threading and the CPU time per task rises by about 40\,\%. In the controlled comparison,
  the same number of executors runs 13\,\% faster on 32 nodes than on 24 nodes, where they share nodes.
  \item The Spark configuration is sound: CPU-bound, 83--95\,\% utilisation, 2--4\,\% GC, no spill and
  low skew. Changing the number of partitions ($\times 3$ vs $\times 6$) has no measurable effect.
  \item The best price/performance is one executor per node. Two executors per node still cut the wall
  time by 16--30\,\% if time matters more than cost.
  \item Scaling does not change the prediction quality (mIoU about 0.38--0.39, OA about 0.84). The
  quality is limited by the capacity of the model and by class imbalance, not by the amount of training data.
\end{enumerate}

\end{document}
